\renewcommand{\headline}{Penta$\cdot$gioco - Italiano}
\renewcommand{\tocent}{Italiano}
\renewcommand{\tocline}{Regole in italiano}
\renewcommand{\translator}{Bozza di Claude, in attesa di revisione}
\renewcommand{\general}{
\lettrine[lines=3]{U}{n} gioco che si può spiegare in un minuto ma che rimane avvincente per anni.

In due si gioca in appena 20 minuti. In tre o quattro si può arrivare fino a 90 minuti.

Non servono dadi. Adatto a tutti dai 5 anni in su.

Nella scatola: 4$\times$5 pedine colorate, 5 blocchi neri e 5 grigi, il tabellone.
}

\renewcommand{\choosext}{Scegli le tue pedine}
\renewcommand{\choosex}{
Ogni giocatore ha cinque pedine. Nella scatola c'è una squadra di pedine per ogni giocatore.

Hai una pedina blu, una rossa, una bianca, una verde e una gialla.

Iniziano ai cinque angoli del tabellone corrispondenti al loro colore.

Vogliono raggiungere le grandi fermate al centro, dove escono dal gioco.
}

\renewcommand{\setupt}{Preparazione}
\renewcommand{\setup}{
Metti le tue pedine sui grandi angoli dell'anello corrispondenti al loro colore: la pedina bianca nell'angolo bianco dell'anello, la pedina blu nell'angolo blu dell'anello, e così via.

Metti i blocchi neri sui cinque incroci al centro del tabellone.

Lascia i blocchi grigi al centro per dopo.
}

\renewcommand{\objectivet}{Obiettivo}
\renewcommand{\objective}{
Le pedine bianche vogliono andare all'incrocio bianco al centro, quelle blu all'incrocio blu, e così via. La destinazione è sempre la grande fermata colorata al centro del pentagono, opposta al punto di partenza.

Sii il primo a spostare \emph{tre} pedine nelle loro destinazioni per vincere.
}

\renewcommand{\rulest}{Le regole}
\renewcommand{\rules}{
Sposta una qualsiasi delle tue pedine sulla stella o sull'anello, in qualsiasi direzione, per quanto vuoi.

Puoi girare in qualsiasi angolo o incrocio libero senza fermarti.

Non puoi mai saltare\textemdash né sopra i blocchi, né sopra altre pedine.

\myskip

Ma puoi muoverti \emph{su} una fermata occupata:

\myskip

Se così facendo catturi un blocco nero, mettilo su una fermata libera a tua scelta.

Se così facendo ti sposti su una fermata con un'altra pedina, scambia la posizione con essa.

\myskip

Puoi scambiare le posizioni di due delle tue pedine.

Se ti sposti su una fermata con più pedine, scegline una con cui scambiare.

\myskip

Non fare due volte di seguito la stessa identica mossa.

\myskip

Una pedina che ha raggiunto la sua destinazione viene tolta dal gioco. Mettila al centro del tabellone. Poi prendi uno dei blocchi grigi e mettilo su una fermata libera a tua scelta.

Se catturi un blocco grigio in questo modo, toglilo di nuovo dal tabellone.

Se finisci i blocchi grigi, sposta uno di quelli già in gioco.

\myskip

Se un altro giocatore ti ha spostato sulla tua meta, devi uscire non appena arriva il tuo turno.

\myskip

L'ultimo giro si gioca fino alla fine.
}
